\section{Conclusion and Future Directions}
\label{sec:conclusion}

In this work, we introduced \textbf{ChemMerge}, a training-free continuous-time paradigm that redefines dynamic model ensembling on edge devices. By challenging the traditional assumption of layer-wise stateless routing, we modeled representation flow through network blocks as a multi-component chemical reactor governed by non-equilibrium reaction kinetics. Through our three components---Catalytic Zero-Shot Alignment (C-ZCA), Non-Equilibrium Kinetic Routing (NEKR), and Catalytic Activation Blending (CAB)---we established physical continuity and representation inertia across sequential layers. 

Our evaluations demonstrate that ChemMerge achieves absolute robustness under extreme task mixtures, recovering \textbf{98.81\%} of the Expert Ceiling and outperforming stateless nearest-centroid baselines by up to \textbf{+8.22\%} while remaining statistically comparable in accuracy to other state-of-the-art activation blending techniques. ChemMerge remains completely immune to both Heterogeneity Collapse and Vectorization Collapse while executing in a single parallel forward pass with constant $O(1)$ latency.

\subsection{Real-World Generalization and Scaling Verification}
We have successfully bridged the gap between synthetic coordinate simulations and real-world deep neural representations through two extensive empirical verification studies:
\begin{itemize}
    \item \textbf{High Expert Density Scaling:} In Section~\ref{sec:expert_scaling}, we scaled the expert count up to $K=16$ experts. We demonstrated that while stateless routing collapses due to high-frequency weight oscillations, ChemMerge maintains robust ensembling accuracy ($74.20\%$, outperforming SPS-ZCA by $+24.10\%$). Profoundly, we found that because ChemMerge's equations are fully vectorized, its parallelized routing solver executes in only $19.9\text{ms}$ ($42.1\%$ faster than SABLE and $49.4\%$ faster than SPS-ZCA), proving outstanding efficiency.
    \item \textbf{Real-World Vision Transformer Generalization:} In Section~\ref{sec:real_world_vit}, we deployed ChemMerge on a pre-trained Vision Transformer ($\text{ViT-B/16}$) model across a heterogeneous shape-classification stream. Using the Multi-Centroid Mode, ChemMerge successfully achieved a high routing accuracy of $93.20\%$ while reducing layer-to-layer ensembling weight jitter by up to $9.9\times$ compared to stateless routing (and $2.15\times$ compared to SABLE under equivalent sensitivity), confirming that continuous kinetics act as a highly effective, robust low-pass filter on actual high-dimensional manifolds.
\end{itemize}

\subsection{Future Research Horizons}
\label{sec:future_horizons}
Our non-equilibrium chemical ensembling framework opens up several exciting, high-impact avenues for future inquiry:
\begin{itemize}
    \item \textbf{Transitioning to Real-World Adapter Ensembling on Standard Benchmarks:} 
    In this work, we verified the mathematical properties of our kinetics routing on real pre-trained Vision Transformers using a routing-only validation. An immediate and high-priority next step is to evaluate ChemMerge's full ensembling capabilities (specifically executing Catalytic Activation Blending, CAB) across actual trained adapters. We present a concrete, five-step roadmap for executing this transition:
    \begin{enumerate}
        \item \textbf{Expert Adapter Training Phase:} First, we train task-specific parameter-efficient adapters (such as LoRA adapters or adapters fine-tuned on MLP layers) for multiple downstream tasks. For vision, this includes the 19 diverse datasets of the Visual Task Adaptation Benchmark (VTAB-1k) \cite{zhai19vtab} or DomainNet's distinct domains. For NLP, we train LoRAs on the 8 classification tasks of the GLUE benchmark \cite{wang18glue}.
        \item \textbf{Offline Calibration Phase:} Next, we extract activation representations from intermediate layers of the shared pre-trained backbone across a tiny calibration set of 16--32 samples per task. We compute and store the layer-wise centroid representations $\mu_k^{(l)}$ for each downstream task expert.
        \item \textbf{Dynamic Serving Pipeline Integration:} During test-time serving, we receive a heterogeneous streaming sequence of unlabelled samples. In the forward pass, at each layer $l$, we compute the cosine similarity between the current hidden activation $h_b^{(l-1)}$ and the layer's centroids $\mu_k^{(l-1)}$. We feed these similarities into the temperature-scaled Arrhenius rate equations (with $\tau = 0.01$) to compute the catalytic reaction rates $k_k^{(l)}$.
        \item \textbf{ODE Kinetics Evolution:} We update the concentration vector $C_k^{(l)}$ using the exact Exponential Integrator (with $\Delta t = 1.5$ and $k_{\text{decay}} = 0.3$). The final blending coefficients $\alpha_k^{(l)}$ are computed via the Law of Mass Action (Eq.~\ref{eq:mass_action}).
        \item \textbf{Parallel Activation Blending (CAB):} Inside the parallel forward pass, we load and execute the specialized LoRAs in parallel, multiplying their output activations by their respective blending coefficients $\alpha_k^{(l)}$ and adding them to the base representation. This end-to-end serving setup will be evaluated under diverse streaming heterogeneity conditions, demonstrating ChemMerge's ability to avoid parameter interference, maintain optimal routing, and suppress representation jitter without requiring any stateful queueing or batch buffering.
    \end{enumerate}
    
    \item \textbf{Scaling to Autoregressive Large Language Models (LLMs):} 
    Applying ChemMerge to massive autoregressive language and multi-modal foundation models (e.g., LLaMA or GPT-series) fine-tuned with task-specific adapters represents a major frontier. 
    This extension would explore ensembling adapters dynamically on a token-by-token or prompt-by-prompt heterogeneous context. 
    Specifically, the concentration state $C_k$ can propagate in two orthogonal physical dimensions: (i) \emph{Intra-token depth propagation}, where concentrations evolve across successive decoder layer blocks ($l = 1 \dots L$) within a single token's forward pass, smoothing layer-to-layer transitions, and (ii) \emph{Inter-token temporal propagation}, where the final concentration $C_k^{(L)}$ of token $t-1$ serves as the boundary condition $C_k^{(0)}$ for token $t$ in time. 
    This creates a dual-axis continuous reactor cascade that stabilizes both the representational depth trajectory and the generation sequence trajectory, preventing multi-turn context collapse and stabilizing generation trajectories under severe conversational topic shifts. 
    To optimize the trade-off between physical transition lag (topic shifts) and routing stability in this temporal setting, we propose \emph{dynamic, similarity-driven temporal memory gating}. 
    By computing the change in representation (or surprise/perplexity) between successive tokens, we can detect discrete topic shifts and instantly trigger a \emph{thermal reset} (setting concentrations back to uniform $1/K$), bypassing transition lag completely on topic boundaries while maintaining smooth ensembling within cohesive sequences.
    \item \textbf{Implicit and Exponential Integration Schemes:} In this work, we successfully derived, implemented, and empirically evaluated an exact **Exponential Integration scheme** (Eq.~\ref{eq:exponential_integrator}) as a mathematically rigorous alternative to Explicit Euler discretization. Because Eq.~\ref{eq:exponential_integrator} represents a strict convex combination of the previous concentration $C_k^{(l-1)} \in [0, 1]$ and the steady-state equilibrium $C_k^* \in [0, 1]$, the updated concentration $C_k^{(l)}$ is mathematically guaranteed to remain strictly within $[0, 1]$ for any virtual step size $\Delta t > 0$ without requiring any heuristic projection clipping. Our empirical evaluations in Section~4 confirm that this scheme remains perfectly stable and high-performing even under extremely large step sizes (e.g., $\Delta t = 10.0$), completely outperforming Explicit Euler. Future work can extend this by exploring implicit multi-step physical solvers (e.g., Runge-Kutta or backward-differentiation formulas) to handle even more stiff kinetics dynamically.
    \item \textbf{Complex Chemical Reaction Networks (CRNs):} While this work successfully leveraged first-order reversible kinetics, more sophisticated biochemical phenomena can be explored. Future research could investigate non-linear dynamical architectures such as bimolecular reactions (modeling direct interaction and synergistic coupling between different expert representations) and autocatalytic feedback loops (such as Gray-Scott or Brusselator systems). Autocatalytic networks, where active species catalyze their own production, could be exceptionally powerful for modeling multi-turn conversation context or prompt-level task reinforcement in LLMs: once a task context is detected, the corresponding expert concentration would self-amplify, reinforcing the contextual trace across multiple conversational turns until active back-reaction decay or a discrete topic shift occurs. We could also model model ensembling as a complex Reaction-Diffusion network with spatial diffusion across parallel attention heads.
    \item \textbf{Neuromorphic and Analog Hardware Co-Design:} Because our kinetics are derived from continuous-time physical ODEs, they are uniquely suited for analog neuromorphic processors or biological computing substrates. These platforms can compute the Euler concentration updates natively using physical material kinetics at near-zero power consumption, paving the way for ultra-low-power, self-organizing edge intelligence.
\end{itemize}
Our chemical-biological ensembling framework represents a powerful step toward creating self-organizing, adaptive, and highly robust intelligence.
